\section{Cloven categories over~$\cal{E}$}
\label{VI.7}

\begin{definition}
\label{VI.7.1} Let $\cal{F}$ be a category over~$\cal{E}$. One
calls \emph{cleavage}
of~$\cal{F}$ over~$\cal{E}$ a function which attaches to every
$f\in\Fl(\cal{E})$ an inverse image functor for $f$ in $\cal{F}$,
namely~$f^*$. The cleavage is said to be \emph{normalized}
if $f=\id_S$ implies $f^*=\id_{\cal{F}_S}$. One calls
\emph{cloven category} (\resp \emph{normalized cloven
category})
a category $\cal{F}$ over~$\cal{E}$ equipped with a cleavage
(\resp with a normalized cleavage).
\end{definition}

It is clear that $\cal{F}$ admits a cleavage if and only if
$\cal{F}$ is prefibered over~$\cal{E}$, and then $\cal{F}$
admits a normalized cleavage. The set of cleavages on~$\cal{F}$
is in bijective correspondence with the set of subsets $K$ of
$\Fl(\cal{F})$ satisfying the following conditions:
\begin{enumerate}
\item[a)] The $\alpha\in K$ are cartesian morphisms.
\item[b)] For every morphism $f\colon T\to S$ in $\cal{E}$ and every
$\xi\in\Ob(\cal{F}_S)$, there exists a unique $f$-morphism in $K$, with
target $\xi$.
\end{enumerate}

For the cleavage defined by $K$ to be normalized, it is necessary and
sufficient that $K$ satisfy in addition the condition
\begin{enumerate}
\item[c)] The identity morphisms in $\cal{F}$ belong
to~$K$.
\end{enumerate}

% original p. 171
The morphisms that are elements of~$K$ may be called the
\emph{``transport morphisms''}
for the cleavage under consideration.

The notion of isomorphism of cloven categories over~$\cal{E}$
is clear. More generally, one may define the
morphisms of cloven $\cal{E}$-categories as the functors
of $\cal{E}$-categories $\cal{F}\to \cal{G}$ which send
transport morphisms to transport morphisms. (These are in
particular cartesian functors). In this way the
cloven categories over~$\cal{E}$ are the objects of a
category, the \emph{category of cloven categories
over} $\cal{E}$. The reader will spell out the existence of products,
related to the fact that if a category over~$\cal{E}$ is a product of
categories $\cal{F}_i$ over~$\cal{E}$ each equipped with a cleavage,
then $\cal{F}$ is equipped with a corresponding natural cleavage. One
also leaves to the reader to spell out the notion of change of base
in cloven categories.

We shall denote by $\alpha_f(\xi)$ the canonical morphism
$$
\alpha_f(\xi)\colon f^*(\xi)\to\xi.
$$
It is, as has been said, functorial in $\xi$, \ie one has a functorial
homomorphism
$$
\alpha_f\colon i_Tf^*\to i_S,
$$
where for every $S\in\Ob(\cal{E})$, $i_S$ denotes the inclusion
functor
$$
i_S\colon\cal{F}_S\to \cal{F}
$$

Consider now morphisms
$$
f\colon T\to S\quad\text{ and }\quad g\colon U\to T
$$
in $\cal{E}$, and let $\xi\in\Ob(\cal{F}_S)$, there then exists a
unique $U$-morphism
$$
c_{f,g}(\xi)\colon g^*f^*(\xi)\to (fg)^*(\xi)
$$
making
% original p. 172
the diagram commutative
$$
\xymatrix{
f^*(\xi) \ar[d]_{\alpha_f(\xi)} & & \ar[ll]_{\alpha_g(f^*(\xi))}
g^*(f^*(\xi)) \ar[d]^{c_{f,g}(\xi)} \\
\xi & & \ar[ll]^{\alpha_{fg}(\xi)} (fg)^*(\xi) }
$$
(by virtue of the definition of~$(fg)^*(\xi)$). For variable $\xi$,
this homomorphism is functorial, {i.e.}: one has a homomorphism
$$
c_{f,g}\colon g^*f^*\to (fg)^*
$$
of functors $\cal{F}_S\to\cal{F}_U$. Let us note at once:
\begin{proposition}
\label{VI.7.2} For the cloven category $\cal{F}$
over~$\cal{E}$ to be fibered, it is necessary and sufficient that the $c_{f,g}$
be isomorphisms.
\end{proposition}

One concludes from this, taking for $f$ an isomorphism, for $g$ its inverse,
and considering the isomorphisms $c_{f,g}$ and $c_{g,f}$:
\begin{corollary}
\label{VI.7.3} If $\cal{F}$ is a \emph{fibered}
cloven category over~$\cal{E}$, then for every isomorphism $f\colon T\to
S$ in $\cal{E}$, $f^*$ is an equivalence of categories
$\cal{F}_S\to \cal{F}_T$.
\end{corollary}

\begin{proposition}
\label{VI.7.4} Let $\cal{F}$ be a cloven category over~$\cal{E}$.
One has

$\qquad\qquad A)\qquad
\left\{\begin{array}{ccc}c_{f,\id_T}(\xi)&=&\alpha_{\id_T}(f^*(\xi))\\
c_{\id_S,f}(\xi)&=&f^*(\alpha_{\id_S}(\xi))
\end{array}\right.$

\smallskip
$\qquad\qquad B)\qquad c_{f,gh}(\xi)\cdot
c_{g,h}(f^*(\xi))=c_{fg,h}(\xi)\cdot h^*(c_{f,g}(\xi))$

\noindent(In
% original p. 173
these formulas, $f,g,h$ denote morphisms
$$
V\to U\to T\to S
$$
and $\xi$ an object of~$\cal{F}_S$).
\end{proposition}

The first and second relations, in the case of a normalized
cleavage, take the simpler form
\begin{equation*}
\label{eq:VI.7.A'}\quad A')\qquad {c_{f,\id_T}=\id_{f^*},\quad
c_{\id_S,f}=\id_{f^*}.\qquad\qquad\qquad\qquad\qquad\qquad}
\end{equation*}

As for the third, it is visualized by the
commutativity of the diagram
\begin{equation*}
\tag{$D$}
\label{eq:VI.7.D}
\begin{split}
\xymatrix@C=2.2cm{
h^*g^*f^*(\xi) \ar[r]^-{c_{g,h}(f^*(\xi))} \ar[d]_{h^*(c_{f,g}(\xi))}
& (gh)^*(f^*(\xi)) \ar[d]^{c_{f,gh}(\xi)} \\
h^*(fg)^*(\xi) \ar[r]_-{c_{fg,h}(\xi)} & (fgh)^*(\xi) }
\end{split}
\end{equation*}
In the case of fibered
categories
(where the $c_{f,g}$ are isomorphisms), this commutativity
can be expressed intuitively by the fact that \emph{the successive
use of isomorphisms of the form $c_{f,g}$ does not lead
to ``contradictory identifications''.} One can equally
write this formula without argument $\xi$, by using the
convolution product of homomorphisms of functors:
$$
c_{fg,h}\circ (h^**c_{f,g})=c_{f,gh}\circ(c_{g,h}*f^*).
$$

The proof of the first two formulas~\ref{VI.7.4} is
trivial; let us sketch that of the third. For this,
consider, in addition to the square $(D)$, the square
of homomorphisms:
\begin{equation*}
\tag{$D'$}
\label{eq:VI.7.D'}
\begin{split}
\xymatrix@C=2cm{
g^*f^*(\xi) \ar[r]^-{\alpha_g(f^*(\xi))}
\ar[d]_{c_{f,g}(\xi)} & f^*(\xi) \ar[d]^{\alpha_f(\xi)} \\
(fg)^*(\xi) \ar[r]_-{\alpha_{fg}(\xi)} & \xi }
\end{split}
\end{equation*}
which
% original p. 174
is commutative by definition of
$c_{f,g}(\xi)$. Consider the diagram obtained by joining the
vertices of~$(D)$ to the corresponding vertices of~$(D')$ by the
homomorphisms of the form $\alpha$:
$$
\begin{array}{ll}
\alpha_h(g^*f^*(\xi)),&\alpha_{gh}(f^*(\xi))\\
\alpha_h((fg)^*(\xi)),&\alpha_{fgh}(\xi).
\end{array}
$$
The four lateral faces of the cube thus obtained are likewise
commutative: for the left-hand one, this comes from the fact that the
left-hand column of~$(D)$ is deduced from the left-hand column of~$(D')$
by applying~$h$, and that $\alpha_h$ is a functorial
homomorphism; for the other three, it is nothing other than the
definition of the operations $c$ of the three remaining sides
of~$(D)$. Thus the five faces of the cube other than the
top face are commutative. It follows that the two
$(fgh)$-morphisms $h^*g^*f^*(\xi)\to(fgh)^*(\xi)$ defined by
$(D)$ have the same composite with $\alpha_{fgh}(\xi)\colon
(fgh)^*(\xi)\to\xi$, hence they are equal by definition of
$(fgh)^*$.

Let us restrict ourselves in what follows to \emph{normalized}
cloven categories. Such a category gives rise to the
following objects:
\begin{enumerate}
\item[a)] A map $S\mto \cal{F}_S$ from~$\Ob(\cal{E})$ to
$\Cat$.
\item[b)] A map $f\mto f^*$, associating to every
$f\in\Fl(\cal{E})$, with source $T$ and target $S$,
a functor $f^*\colon \cal{F}_S\to\cal{F}_T$.
\item[c)] A map $(f,g)\mto c_{f,g}$, associating to every
pair of arrows $(f,g)$ of~$\cal{E}$, a functorial
homomorphism $c_{f,g}\colon g^*f^*\to (fg)^*$.
\end{enumerate}

Moreover these data satisfy the conditions expressed in
the formulas $A')$ and~$B)$ given above. (N.B. If one had not
restricted oneself to the case of a normalized cleavage, it would have
been necessary to introduce an additional object, namely a function
$S\mto \alpha_S$ which associates to every object $S$ of~$\cal{E}$ a
functorial homomorphism $\alpha_S\colon (\id_S)^*\to
\id_{\cal{F}_S}$; the condition $A')$ would then be replaced by the
condition $A)$).

% original p. 175
We shall now show how one can reconstruct (up to
unique isomorphism) the normalized cloven
category $\cal{F}$ over~$\cal{E}$ with the help of the preceding
objects.

\section[Cloven category defined by a pseudofunctor]{Cloven category defined by a pseudofunctor
$\cal{E}^\circ\to\Cat$}
\label{VI.8}

Let us call, for brevity, a \emph{pseudofunctor} from~$\cal{E}^\circ$
to $\Cat$ (one should say, a \emph{normalized} pseudofunctor),
a set of data a), b), c) as above, satisfying the
conditions $A')$ and $B)$. In the preceding no., we
associated, to a normalized cloven category
over~$\cal{E}$, a pseudofunctor $\cal{E}^\circ\to\Cat$; here we
shall indicate the inverse construction. We shall leave to the reader the
verification of most of the details, as well as of the fact that
these constructions are indeed ``inverse'' to one another. More
precisely, one should consider the
pseudofunctors $\cal{E}^\circ\to\Cat$ as the objects of a
new category, and show that our constructions
furnish equivalences, quasi-inverse to one another,
between the latter and the category of
cloven categories over~$\cal{E}$, defined in the preceding
no.

One sets
$$
\cal{F}_\circ=\underset{S\in\Ob(\cal{E})}\coprod
\Ob(\cal{F}(S)),
$$
the sum of the sets
$\Ob(\cal{F}(S))$
(N.B. we shall write here $\cal{F}(S)$ and not $\cal{F}_S$ for the value at
the object $S$ of~$\cal{E}$ of the given pseudofunctor, to avoid
confusions of notation in what follows). One thus has an obvious
map:
$$
p_\circ\colon\cal{F}_\circ\to\Ob\cal{E}.
$$
Let
$$
\overline{\xi}
=(S,\xi),\quad \overline{\eta}=(T,\eta)\qquad\qquad
\text{(with $\xi\in\Ob\cal{F}(S)$, $\eta\in\Ob\cal{F}(T)$)}
$$
be two elements of~$\cal{F}_\circ$, and let $f\in\Hom(T,S)$, one
will set
$$
h_f(\overline{\eta},\overline{\xi})=\Hom_{\cal{F}(T)}(\eta,f^*(\xi)).
$$
% original p. 176
If moreover one has a morphism $g\colon U\to T$ in $\cal{E}$, and a
$\zeta\in\Ob\cal{F}(U)$, one defines a map, denoted
$(u,v)\mto u\circ v$:
$$
h_f(\overline{\eta},\overline{\xi})\times
h_g(\overline{\zeta},\overline{\eta})\to
h_{fg}(\overline{\zeta},\overline{\xi}),
$$
\ie a map
$$
\Hom_{\cal{F}(T)}(\eta,f^*(\xi))\times\Hom_{\cal{F}(U)}(\zeta,g^*(\eta))
\to \Hom_{\cal{F}(U)}(\zeta,(fg)^*(\xi)),
$$
by the formula
$$
u\circ v=c_{f,g}(\xi)\cdot g^*(u)\cdot v,
$$
\ie $u\circ v$ is the composite of the sequence
$$
\zeta\lto{u}
g^*(\eta)\lto{g^*(u)}
g^*f^*(\xi)\lto{c_{f,g}(\xi)} (fg)^*(\xi).
$$
On the other hand one will set
$$
h(\overline{\eta},\overline{\xi})=\underset{f\in\Hom(T,S)}\coprod
h_f(\overline{\eta},\overline{\xi}),
$$
and the preceding pairings define
pairings
$$
h(\overline{\eta},\overline{\xi})\times
h(\overline{\zeta},\overline{\eta})\to
h(\overline{\zeta},\overline{\xi}),
$$
while the definition of the $h(\overline{\eta},\overline{\xi})$
implies an obvious map:
$$
p_{\overline{\eta},\overline{\xi}}\colon
h(\overline{\eta},\overline{\xi})\to \Hom(T,S).
$$
This being said, one verifies the following points:
\begin{enumerate}
\item[1)] The composition between elements of the
$h(\overline{\eta},\overline{\xi})$
is \emph{associative}.
\item[2)]
% original p. 177
For every $\overline{\xi}=(\xi,S)$ in $\cal{F}_\circ$,
consider the identity element
of
$$
h_{\id_S}(\overline{\xi},\overline{\xi})=
\Hom_{\cal{F}(S)}(\id_S^*(\xi),\xi)=\Hom_{\cal{F}(S)}(\xi,\xi),
$$
and its image in $h(\overline{\xi},\overline{\xi})$. This object is
a \emph{unit} on the left and on the right for the composition
between elements of the $h(\overline{\eta},\overline{\xi})$.

This already shows that \emph{one obtains a category}
$\cal{F}$, by setting
$$
\Ob\cal{F}=\cal{F}_\circ,\quad \Fl{\cal{F}}=
\underset{\overline{\xi},\overline{\eta}\in\cal{F}_\circ}
\coprod h(\overline{\eta},\overline{\xi}).
$$
(N.B. one cannot simply take for $\Fl\cal{F}$ the
\emph{union} of the sets $h(\overline{\eta},\overline{\xi})$,
for the latter are not necessarily disjoint). Moreover:
\item[3)] The maps $p_\circ\colon \Ob\cal{F}\to\Ob\cal{E}$ and
$p_1=(p_{\overline{\eta},\overline{\xi}})\colon\Fl\cal{F}\to\Fl{\cal{E}}$
define a \emph{functor} $p\colon\cal{F}\to\cal{E}$. In
this way, $\cal{F}$ becomes a category over~$\cal{E}$, and
moreover the obvious map
$h_f(\overline{\eta},\overline{\xi})\to\Hom(\overline{\eta},\overline{\xi})$
induces a \emph{bijection}
$$
h_f(\overline{\eta},\overline{\xi})\isomto
\Hom_f(\overline{\eta},\overline{\xi}).
$$
\item[4)] The obvious maps
$$ \Ob\cal{F}(S)\to\cal{F}_\circ=\Ob\cal{F},\qquad \Fl\cal{F}(S)\to
\Fl\cal{F},
$$
where the second is defined by the obvious
maps
$$
\Hom_{\cal{F}(S)}(\xi,\xi')=
h_{\id_S}(\overline{\xi},\overline{\xi}')\to
\Hom(\overline{\xi},\overline{\xi}')
$$
define an \emph{isomorphism}
$$
i_S \colon \cal{F}(S)\isomto\cal{F}_S.
$$
\item[5)] For every object $\overline{\xi}=(S,\xi)$ of~$\cal{F}$, and
every morphism $f\colon T\to S$ of~$\cal{E}$, consider
% original p. 178
the element $\overline{\eta}=(T,\eta)$ of~$\cal{F}_T$, with
$\eta=f^*(\xi)$, and the element $\alpha_f(\xi)$ of
$\Hom(\overline{\eta},\overline{\xi})$, image of~$\id_{f^*(\xi)}$ by
the morphism
$\Hom_{\cal{F}(T)}(f^*(\xi),f^*(\xi))=h_f(\overline{\eta},\overline{\xi})
\to \Hom_f(\overline{\eta},\overline{\xi})$. \emph{This element
is cartesian, and it is the identity in $\overline{\xi}$ if}
$f=\id_S$, in other words, the set of the $\alpha_f(\xi)$
defines a \emph{normalized cleavage of~$\cal{F}$ over}
$\cal{E}$. Moreover, by construction, one has commutativity in the
diagram of functors
$$
\xymatrix{
\cal{F}(S) \ar[r]^{f^*} \ar[d]_{i_S} & \cal{F}(T) \ar[d]^{i_T} \\
\cal{F}_S \ar[r]^{f^*_{\cal{F}}} & \cal{F}_T }
$$
where $f^*_{\cal{F}}$ is the inverse image functor by $f$, relative
to the cleavage under consideration on~$\cal{F}$. Finally:
\item[6)] the homomorphisms $c_{f,g}$ given with the
pseudofunctor are transformed, by the isomorphisms $i_S$, into
the functorial homomorphisms $c_{f,g}$ associated to the cleavage
of~$\cal{F}$.
\end{enumerate}

We confine ourselves to giving the verification of 1) (which is, if
possible, less trivial than the others). It suffices to prove
the associativity of the composition between the objects of sets of
the form $h_f(\overline{\eta},\overline{\xi})$. Consider therefore
in $\cal{E}$ morphisms
$$
S\lfrom{f} T\lfrom{g} U\lfrom{h} V
$$
and objects
$$
\xi,\ \eta,\ \zeta,\ \tau
$$
in $\cal{F}(S),\cal{F}(T),\cal{F}(U),\cal{F}(V)$, finally
elements
$$
\begin{array}{ccc}
u\in
h_f(\overline{\eta},\overline{\xi})&=&\Hom_{\cal{F}(T)}(\eta,f^*(\xi))\\
v\in
h_g(\overline{\zeta},\overline{\eta})&=&\Hom_{\cal{F}(U)}(\zeta,g^*(\eta))\\
w\in
h_h(\overline{\tau},\overline{\zeta})&=&\Hom_{\cal{F}(V)}(\tau,h^*(\zeta)).
\end{array}
$$
One
% original p. 179
wants to prove the formula
$$
(u\circ v)\circ w = u\circ (v\circ w),
$$
which is an equality in
$\Hom_{\cal{F}(V)}(\tau,(fgh)^*(\xi))$. By virtue of the definitions
the two sides of this equality are obtained by composition
along the upper and lower contours of the diagram
below:
$$
\xymatrix{
\tau \ar[r]^w \ar[drr]_{v\circ w} & h^*(\zeta)
\ar[r]^{h^*(v)} \ar@<-2pt>
`u[rrrrr] `[rrrrr]^{h^*(u\circ v)} [rrrrr] & h^*g^*(\eta)
\ar[rr]^{h^*g^*(u)}
\ar[d]_{c_{g,h}(\eta)} & & h^*g^*f^*(\xi) \ar[rr]^{h^*(c_{f,g}(\xi))}
\ar[d]^{c_{g,h}(f^*(\xi))} & & h^*(fg)^*(\xi) \ar[d]^{c_{fg,h}(\xi)} \\
& & (gh)^*(\eta) \ar[rr]^{(gh)^*(u)} & & (gh)^*f^*(\xi)
\ar[rr]^{c_{f,gh}(\xi)} & & (fgh)^*(\xi) }
$$
Now the middle square is commutative because $c_{g,h}$ is a
functorial homomorphism, and the right-hand square is commutative by
virtue of the condition~$B)$ for a pseudofunctor. Whence the
announced result.

Of course, it remains to specify, when the pseudofunctor
under consideration already comes from a normalized cloven
category $\cal{F}'$ over~$\cal{E}$, how one obtains a
natural isomorphism between $\cal{F}'$ and $\cal{F}$. We leave the
details to the reader.

We also leave to the reader to interpret, in terms of
pseudofunctors, the notion of inverse image of a cloven
category $\cal{F}$ over~$\cal{E}$ by a change of base functor
$\cal{E}' \to\cal{E}$.

\section[Example]{Example: cloven category defined by a functor
$\cal{E}^\circ\to\Cat$; split categories over~$\cal{E}$}
\label{VI.9}
Suppose one has a functor
$$
\varphi\colon\cal{E}^\circ\to\Cat,
$$
it then defines a pseudofunctor by setting
$$
\cal{F}(S)=\varphi(S),
\quad f^*=\varphi(f),\quad c_{f,g}=\id_{(fg)^*}
$$
% original p. 180
Thus the construction of the preceding no. gives us a
category $\cal{F}$ cloven over~$\cal{E}$, said to be associated to the
functor $\varphi$. For a cloven category over~$\cal{E}$
to be isomorphic to a cloven category defined by a
functor $\varphi:\cal{E}^\circ\to\Cat$, it is necessary and sufficient,
manifestly, that it satisfy the conditions:
$$
(fg)^*=g^*f^*,\quad c_{f,g}=\id_{(fg)^*}\quoi.
$$
In terms of the set $K$ of transport morphisms, this also simply
means that \emph{the composite of two transport
morphisms is a transport morphism}. A cleavage of a
category $\cal{F}$ over~$\cal{E}$ satisfying the preceding
condition is called a \emph{splitting}
of~$\cal{F}$ over~$\cal{E}$, and a category $\cal{F}$
over~$\cal{E}$ equipped with a splitting is called a
\emph{split category over~$\cal{E}$}.
It is thus a particular case of the notion of
cloven category. The category of split categories
over~$\cal{E}$ is thus equivalent to
$\SheafHom(\cal{E}^\circ,\Cat)$. Note that a
split category over~$\cal{E}$ is a fortiori a cloven category
over~$\cal{E}$.

If $\cal{F}$ is a fibered category over~$\cal{E}$, there
does not always exist a splitting on~$\cal{F}$. Suppose for example
that $\Ob\cal{E}$ and $\Ob\cal{F}$ are reduced to a
single element, and that the set of endomorphisms of the said element is a
group $E$ \resp $F$, so that the projection functor $p$ is
given by a homomorphism of groups $p\colon F\to E$, surjective
since $p$ is fibrant. One then verifies at once that
the set of cleavages of~$\cal{F}$ over~$\cal{E}$ is in
bijective correspondence with the set of maps $s\colon
E\to F$ such that $ps=\id_E$ (\ie the set of ``systems of
representatives'' for the classes mod the subgroup $G$ kernel of
the surjective homomorphism $p\colon F\to E$). A cleavage is a
splitting if and only if $s$ is a homomorphism of groups. To say
that there exists a splitting therefore means that the group extension $F$
of~$E$ by $G$ is trivial, which is expressed, when $G$ is
commutative, by the vanishing of a certain cohomology class
in $\H^2(E,G)$ (where $G$ is regarded as a group
on which $E$ acts).

Suppose however that $\cal{F}$ is a fibered category
over~$\cal{E}$ such that the
% original p. 181
$\cal{F}_S$ are \emph{rigid} categories, \ie the group
of automorphisms of every object of~$\cal{F}_S$ is reduced to
the identity. It is then easy to prove that $\cal{F}$ admits a
splitting over~$\cal{E}$. Indeed, one first observes that the question
of existence of a splitting is not modified if one replaces
$\cal{F}$ by an $\cal{E}$-equivalent category, which in the present
instance brings us back to the case where the $\cal{F}_S$ are
rigid \emph{and reduced} categories (\ie two isomorphic
objects in $\cal{F}_S$ are identical). But if $G$ is a
rigid and reduced category, every isomorphism of two
functors $H\to G$ (where $H$ is an arbitrary category) is
an identity. It follows that if $\cal{F}$ is a
fibered category over~$\cal{E}$, such that the fiber-categories are
rigid and reduced, then there exists a \emph{unique} cleavage of
$\cal{F}$ over~$\cal{E}$, which is necessarily a splitting. Thus
$\cal{F}$ is isomorphic to the category defined by a
functor $\varphi\colon \cal{E}^\circ\to \Cat$, such that the
$\varphi(S)$ are rigid and discrete categories, and the
functor $\varphi$ is defined up to isomorphism.

% ---- en-chunk4.tex ----
